\section{Hotword Conditioning}
\label{sec:hotwords}

\subsection{Setup}

We construct CommonVoice 17 hotword test sets by sampling, for each
utterance in the standard CV17 test split, $k$ ``hotword'' terms drawn
from the union of (i)~rare named entities present in the reference
transcript and (ii)~distractor terms drawn from the global rare-word
distribution that are \emph{not} present in the transcript. We sweep
$k \in \{0, 5, 10, 15, 20\}$ and report CER for Mandarin and WER for
English. The $k=0$ row is identical to the corresponding row of
Tables~\ref{tab:asr_zh} and \ref{tab:asr_en} and serves as the
no-hotword reference.
% TODO: confirm the exact distractor-to-positive ratio used to build the
% test sets, and whether the same ratio is used during Stage~3 training.

\subsection{Results}

\begin{table}[t]
  \caption{Effect of in-prompt hotwords on CommonVoice 17. CER (\%) for
    Mandarin and WER (\%) for English. Lower is better. $k$ is the
    number of hotwords supplied through the prompt.}
  \label{tab:hotwords}
  \centering
  \small
  \begin{tabular}{lrccccc}
    \toprule
    & & $k=0$ & $k=5$ & $k=10$ & $k=15$ & $k=20$ \\
    \midrule
    CV17 zh hotword test  & $n=10{,}647$ & 6.57 & \textbf{2.83} & 2.85 & 2.89 & 2.91 \\
    CV17 en hotword test  & $n=16{,}294$ & 9.65 & \textbf{4.89} & 4.99 & 5.03 & 5.09 \\
    \bottomrule
  \end{tabular}
\end{table}

\subsection{Analysis}

Two effects are visible in Table~\ref{tab:hotwords}:

\paragraph{Large absolute gains from a small list.}
Five hotwords already deliver a $3.74$ absolute CER reduction on
Mandarin (from $6.57$\% to $2.83$\%) and a $4.76$ absolute WER reduction
on English (from $9.65$\% to $4.89$\%). This is consistent with the
intuition that the rare-entity tail dominates the residual error of an
otherwise-well-trained ASR model, and that an LLM is well placed to
``snap'' to the suggested term when it is acoustically supported.

\paragraph{Diminishing -- and slightly inverted -- returns past five.}
Going from five to twenty hotwords has a small but consistent negative
effect (roughly $+0.08$ CER on Mandarin, $+0.20$ WER on English).
% TODO: verify that this monotone pattern persists when we control for
% the number of true-positive terms in the list -- the current sweep
% mixes positives and distractors and the distractor count grows with
% k.
We attribute this in part to false-positive triggering on distractors;
Section~\ref{sec:analysis} returns to this with a per-utterance
breakdown of insertion versus substitution errors.

\paragraph{Implication for deployment.}
For latency- and quality-sensitive deployments, supplying around five
likely hotwords per utterance (rather than dumping a full domain
glossary) appears to be a strong default for AmphionASR.
